% year: תש"פ
% type: מבחן
% semester: א
% moed: ב
% part: ב
% number: 2
% points: 9
% categories: func-analysis, polar-parametric, multiple-choice
% source: exams/מבחנים/תשפ/מועד ב תשפ.pdf

\begin{question}
הגרף הנתון הוא גרף בקואורדינטות פולריות של הפונקציה:

(בשרטוט: שתי לולאות סגורות ושוות, בצורת עיגול, אחת מעל ציר $x$ ואחת מתחתיו, הנוגעות זו בזו בראשית הצירים; הלולאות סימטריות סביב ציר $y$ ונמתחות לאורכו.)
\begin{enumerate}
  \item $r(\phi)=\sin^2\phi$
  \item $r(\phi)=1+\cos\phi$
  \item $r(\phi)=1-\cos\phi$
  \item $r(\phi)=\cos^2\phi$
  \item $r(\phi)=\cos\phi$
\end{enumerate}
\end{question}

\begin{hint}
בדקו עבור כל פונקציה באילו זוויות $r=0$ (שם העקומה עוברת בראשית) ובאילו זוויות $r$ מקסימלי (הנקודות הרחוקות ביותר מהראשית).
\end{hint}

\begin{solution}
\textbf{התשובה הנכונה: א.}

בגרף העקומה עוברת בראשית בכיוון ציר $x$ (שם שתי הלולאות נוגעות), והנקודות הרחוקות ביותר מהראשית נמצאות על ציר $y$, למעלה ולמטה, באותו מרחק. כמו כן יש שתי לולאות.

\begin{itemize}
  \item $r=\sin^2\phi$: מתקיים $r\ge0$ לכל $\phi$, $r=0$ בדיוק עבור $\phi=0,\pi$ (כיוון ציר $x$), ו-$r=1$ מקסימלי עבור $\phi=\frac{\pi}{2},\frac{3\pi}{2}$, כלומר בנקודות $(0,1)$ ו-$(0,-1)$. כש-$\phi$ עובר מ-$0$ ל-$\pi$ נוצרת לולאה בחצי העליון, וכש-$\phi$ עובר מ-$\pi$ ל-$2\pi$ נוצרת לולאה בחצי התחתון. זה בדיוק הגרף הנתון.
  \item $r=1\pm\cos\phi$: קרדיואידות --- לולאה אחת בלבד (עם "שקע" בראשית), סימטרית סביב ציר $x$. לא מתאים.
  \item $r=\cos^2\phi$: שתי לולאות, אבל לאורך ציר $x$ (מקסימום ב-$\phi=0,\pi$). לא מתאים.
  \item $r=\cos\phi$: מעגל אחד, $x^2+y^2=x$, שמרכזו $(\frac12,0)$. לא מתאים.
\end{itemize}
\end{solution}
